\begin{table}[H]
\centering
\caption{Baseline Differences Between Recruited and Non-Recruited Baseline Respondents}
\label{tab:recruitment_selection}
\footnotesize
\begin{threeparttable}
\begin{tabular}{lccccc}
\toprule
Variable & Recruited & Not recruited & Diff & p-value & N \\
\midrule
Baseline attention check passed & 0.999 & 0.869 & 0.130*** & <0.001 & 12,215 \\
Baseline straightlining flag & 0.001 & 0.066 & -0.065*** & <0.001 & 12,215 \\
Study-abroad plan (1=yes, 4=no) & 1.587 & 2.017 & -0.429*** & <0.001 & 12,215 \\
Foreign-news exposure (30 days) & 2.442 & 2.329 & 0.113*** & <0.001 & 12,215 \\
Trust in foreign media & 2.865 & 2.738 & 0.127*** & <0.001 & 12,215 \\
Regime support & 0.164 & 0.376 & -0.212*** & <0.001 & 12,215 \\
Nationalism & -0.016 & 0.146 & -0.161*** & <0.001 & 12,215 \\
Censorship support & -0.092 & 0.100 & -0.192*** & <0.001 & 12,215 \\
China-West thermometer gap & 17.227 & 19.989 & -2.762*** & <0.001 & 12,215 \\
Exam score & 68.037 & 67.609 & 0.428 & 0.242 & 9,558 \\
Female & 0.551 & 0.549 & 0.002 & 0.845 & 12,215 \\
\bottomrule
\end{tabular}
\begin{tablenotes}[flushleft]
\footnotesize
\item Note: Entries compare baseline respondents who entered the randomized study sample (`recruited = 1`) with baseline respondents who did not. Differences are estimated from bivariate OLS regressions with HC2 standard errors. Recruitment was partially based on baseline quality checks and study-abroad interest, so this table is informative about sample selection and external validity rather than experimental balance.
\end{tablenotes}
\end{threeparttable}
\end{table}
